\section{Bayes-factor computation: explicit proofs}
\label{sec:annex-p15-bayes-factor-computation}

\begin{annexitem}{annex:p15.laplace-relative-error}{thm:p15-laplace-relative-error}{Localization, odd-term cancellation and the tail bound}
\begin{proof}
All constants below can depend on the fixed bounds and \(d\).
The local amplitude is bounded by \(a_++A\rho\). Local quadratic
domination and the tail hypothesis make the integral finite. Continuity
and \(a_n(\theta_n)>0\) make it positive.

Put \(u=\sqrt n(\theta-\theta_n)\) and \(b_n=n^{1/12}\).
For large \(n\), the ball \(\|u\|\leq b_n\) lies in the permitted
neighborhood. Taylor's theorem with its bounded fourth derivative gives
\[
 n(h_n(\theta_n+u/\sqrt n)-h_n(\theta_n))
 =-\tfrac12u^TJ_nu+n^{-1/2}T_n(u)+R_n(u),
\]
where \(T_n(u)=D^3h_n(\theta_n)[u,u,u]/6\) is odd,
\(|T_n(u)|\leq A\|u\|^3/6\), and
\(|R_n(u)|\leq A\|u\|^4/(24n)\).
Similarly
\[
 a_n(\theta_n+u/\sqrt n)=a_n(\theta_n)
       +n^{-1/2}Da_n(\theta_n)[u]+O(n^{-1}\|u\|^2).
\]
On this ball the absolute exponent perturbation is at most
\(C(n^{-1/2}b_n^3+n^{-1}b_n^4)\), tending to zero.
The elementary Taylor estimate
\(|e^t-1-t|\leq e\,t^2/2\) for \(|t|\leq1\) follows from the
second-derivative remainder. It shows that the product of the amplitude
and perturbed exponential, after removing \(e^{-u^TJ_nu/2}\), equals
\[
 a_n(\theta_n)+n^{-1/2}
 \{a_n(\theta_n)T_n(u)+Da_n(\theta_n)[u]\}
             +O\{n^{-1}(1+\|u\|^{10})\}.
\]
The constants are uniform. The order-ten polynomial is a convenient bound
for all products of the displayed remainders, not an omitted smoothness
assumption.

The first-order term is odd. Its integral against the even Gaussian on
the centered ball is exactly zero. Every fixed polynomial moment of
\(e^{-c\|u\|^2/2}\) is finite: a polynomial times
\(e^{-c\|u\|^2/4}\) is bounded, leaving an integrable product Gaussian.
Also, for fixed polynomial \(P\),
\[
 \int_{\|u\|>b}|P(u)|e^{-c\|u\|^2/2}du
 \leq e^{-cb^2/4}\int |P(u)|e^{-c\|u\|^2/4}du.
\]
Thus replacing the small ball's Gaussian integral by the full integral
changes it by \(o(n^{-1})\).

For clarity the full Gaussian integral is
\((2\pi)^{d/2}/\sqrt{\det J_n}\). This follows without an unstated
multivariate Jacobian: partition a positive definite matrix as
\(\left(\begin{smallmatrix}a&b^T\\b&C\end{smallmatrix}\right)\).
Complete the square in the first coordinate. Integration gives
\(\sqrt{2\pi/a}\), leaving the positive definite Schur complement
\(C-bb^T/a\), with determinant \(\det J/a\). Induction and the
one-dimensional Gaussian integral prove the formula. Tonelli justifies
the integrations.

Between \(b_n\) and \(\rho\sqrt n\), local quadratic domination bounds
the rescaled integrand by
\((a_++A\rho)e^{-c_0\|u\|^2}\); its integral is \(o(n^{-1})\).
Outside the fixed local ball, the rescaled normalized integral is at most
\(n^{d/2}A_Te^{-\eta n}=o(n^{-1})\). Combining these regions gives
\[
 n^{d/2}e^{-nh_n(\theta_n)}I_n
 =a_n(\theta_n)(2\pi)^{d/2}(\det J_n)^{-1/2}+O(n^{-1}).
\]
The main term is bounded below by
\(a_-(2\pi)^{d/2}C^{-d/2}>0\), proving relative accuracy.
For two models write \(I_k=\widehat I_k(1+\varepsilon_k)\),
\(\varepsilon_k=O(n^{-1})\). Dividing yields the factor assertion.
For large \(n\), \(|\varepsilon_k|\leq1/2\), and
\(|\log(1+\varepsilon_k)|\leq2|\varepsilon_k|\), proving the log assertion.
\end{proof}
\end{annexitem}

\begin{annexitem}{annex:p15.laplace-mode-variants}{thm:p15-laplace-mode-variants}{Substituting the two exponent/amplitude choices}
\begin{proof}
For the posterior choice, \(e^{nh_n}=L_n\pi\), amplitude one, and
\(J_n=H_n/n\). Thus
\(n^{-d/2}(\det J_n)^{-1/2}=(\det H_n)^{-1/2}\).
Substitution into the preceding theorem gives the posterior-mode formula
and its relative error. For the likelihood choice,
\(e^{nh_n}=L_n\), amplitude \(\pi\), and \(J_n=F_n/n\).
The same determinant identity gives the MLE formula. Positivity,
smoothness, local domination and global tails must hold for the chosen
decomposition; they are not transferred from one decomposition by assertion.
\end{proof}
\end{annexitem}

\begin{annexitem}{annex:p15.laplace-curvature-replacement}{thm:p15-laplace-curvature-replacement}{A quantitative determinant comparison}
\begin{proof}
Let \(E=J_n^{-1/2}(\widetilde J_n-J_n)J_n^{-1/2}\).
Its operator norm is at most \(\delta_n/c<1/2\). Its real eigenvalues
\(\epsilon_j\) therefore lie in that interval, and
\(\det\widetilde J_n/\det J_n=\prod_j(1+\epsilon_j)>0\).
Since \(|\log(1+t)|\leq2|t|\) for \(|t|<1/2\),
\[
 \left|\tfrac12\log\frac{\det J_n}{\det\widetilde J_n}\right|
 \leq\sum_j|\epsilon_j|\leq d\delta_n/c.
\]
Only this determinant changes in the specified approximation, proving the
bound. Exponentiating an \(O(\delta_n)\) log difference gives
\(1+O(\delta_n)\), and multiplication by the already proved
\(1+O(n^{-1})\) approximation yields the stated orders.
\end{proof}
\end{annexitem}

\begin{annexitem}{annex:p15.bic-leading-order}{thm:p15-bic-leading-order}{The exact omitted order in the logarithm}
\begin{proof}
Take logarithms of the MLE approximation and use
\(\det F_{k,n}=n^{d_k}\det J_{k,n}\) and the relative-error theorem.
This gives the displayed expression for \(\log I_{k,n}\). Uniform
eigenvalue bounds bound \(\log\det J_{k,n}\); the upper/lower prior
ordinate bounds bound its logarithm. Hence \(c_{k,n}=O(1)\).
Subtract the two expansions to obtain the BIC relation. Multiplying
the unshortened relation by two gives coefficient two on both prior-log
terms, as displayed in the body.

Let \(S_n=-\tfrac12(\mathrm{BIC}_1-\mathrm{BIC}_0)\).
Then \(\log B_{10}=S_n+c_{1,n}-c_{0,n}+O(n^{-1})\), so
\(e^{S_n}/B_{10}\to e^{-c_*}\) when the constant difference converges.
If \(|\log B_{10}|\to\infty\), dividing the bounded additive difference
by \(\log B_{10}\) gives zero. A bounded denominator supplies no such
conclusion. Parameter-free positive likelihoods need no approximation
and contribute no determinant or parameter-prior term.
\end{proof}
\end{annexitem}

\begin{annexitem}{annex:p15.laplace-normal-example}{thm:p15-laplace-normal-example}{Exact normal evidence and failed hypotheses}
\begin{proof}
Write \(s^2=\sigma^2/n\). The completed-square calculation in
\cref{thm:p15-normal-point-null-factor} gives
\(I_n=\phi_{\mu,s^2+\tau^2}(x_n)\) and a product likelihood/prior
whose log is quadratic, with mode
\[
 \widetilde\theta_n=\frac{\tau^2x_n+s^2\mu}{s^2+\tau^2},
 \qquad H_n=1/s^2+1/\tau^2.
\]
Thus \(L_n(\theta)\pi(\theta)\) is exactly \(I_n\) times the normal
density with this mean and precision. Its value at the mode times
\(\sqrt{2\pi/H_n}\) is exactly \(I_n\).
The MLE is \(x_n\), with \(F_n=1/s^2\), and
\(\sqrt{2\pi s^2}L_n(x_n)\pi(x_n)=\pi(x_n)\).
Dividing this normal density by the exact convolution gives the ratio
in the theorem.

For bounded \(x_n\), the prior ordinate at the MLE is bounded away from
zero. Gaussian prior first and second derivatives are uniformly bounded.
Both exponents have zero third and fourth derivatives and curvatures
at least \(1/\sigma^2\) after division by \(n\).
Outside a fixed radius \(\rho\), their normalized Gaussian integrands
are bounded by a constant times
\(e^{-n\rho^2/(4\sigma^2)}e^{-n v^2/(4\sigma^2)}\).
The latter factor integrates to a bounded value for \(n\geq1\), proving
the required exponential-tail bound; the prior amplitude is bounded in
the MLE case. This verifies every Laplace hypothesis.
The logarithm of the explicit MLE ratio is \(O(s^2)=O(n^{-1})\)
for bounded \(x_n\), giving its claimed relative rate.

At \(x_n=\mu=0,\sigma^2=1,\tau^2=4\), divide
\(I_n=1/\sqrt{2\pi(4+1/n)}\) by the null density
\(\sqrt{n/(2\pi)}\). This gives \(1/\sqrt{1+4n}\).
The MLE approximation uses prior ordinate \(1/(2\sqrt{2\pi})\),
giving \(1/(2\sqrt n)\). The maximized alternative likelihood equals
the null likelihood here, so the one-parameter BIC penalty gives
\(1/\sqrt n\). Its ratio to the exact factor tends to two.

For the boundary case, substitution \(u=\sqrt n\theta\) yields
\[
 I_n=n^{-1/2}\int_0^{\sqrt n}e^{-u^2/2}du
 \sim\tfrac12\sqrt{2\pi/n}.
\]
There is no two-sided interior ball at zero. For the two-mode example,
the prior density is \(1/3\) on a total length three; symmetry gives
\[
 I_n=\frac23\int_{1/2}^{2}e^{-n(\theta-1)^2/2}d\theta
 =\frac2{3\sqrt n}\int_{-\sqrt n/2}^{\sqrt n}e^{-u^2/2}du
 \sim\frac23\sqrt{2\pi/n}.
\]
One mode's local contribution is only one half of this, and the other
mode violates the exponentially negligible outside-tail condition.
For the quartic case, symmetry and \(u=n^{1/4}\theta\) give
\[
 I_n=n^{-1/4}\int_0^{n^{1/4}}e^{-u^4}du.
\]
The limiting integral is positive and finite: above one,
\(e^{-u^4}\leq e^{-u}\). The Hessian at zero vanishes, so the positive
curvature premise fails. Finally, each exponential for one trial is in
\((0,1]\); taking it as a Bernoulli success probability and observing
\(n\) successes constructs the stated likelihoods under proper priors.
\end{proof}
\end{annexitem}

\begin{annexitem}{annex:p15.direct-evidence-mc}{thm:p15-direct-evidence-mc}{The integrable prior average}
\begin{proof}
Since \(L\geq0\) and \(\int Ld\Pi=I<\infty\), the sampled likelihoods
are integrable identically distributed variables. Linearity gives
\(E\widehat I_M=I\). The i.i.d. sampling meets the stronger-than-needed
independence hypothesis of \cref{thm:p15-etemadi-strong-law}, so the
average converges almost surely to \(I\). If the second moment is finite,
independence eliminates cross-covariances, giving
\(\operatorname{Var}(\widehat I_M)=\operatorname{Var}_\Pi(L)/M\).
No finite second moment was used for almost-sure convergence.
\end{proof}
\end{annexitem}

\begin{annexitem}{annex:p15.importance-evidence}{thm:p15-importance-evidence}{Domination and two different denominators}
\begin{proof}
Radon--Nikodym change of measure gives
\(\int V\,dQ=T(\Theta)=I\). Set \(V=0\) on its null exceptional set
to make it finite everywhere. The sampled variables are integrable
i.i.d.; the same strong law and variance calculation prove the ordinary
importance assertions.

If \(\Pi\ll Q\), choose \(W=d\Pi/dQ\), likewise finite everywhere
after a null-set assignment. Its expectation under \(Q\) is one, and
change of measure gives \(E_Q[WL]=I\), with \(WL\geq0\) integrable.
The two strong laws hold on a common probability-one event:
\[
 M^{-1}\sum_iW_i\to1,\qquad M^{-1}\sum_iW_iL_i\to I.
\]
The denominator is eventually positive and their ratio tends to \(I\).
If \(\pi,q\) are densities, \(\pi/q\) on \(q>0\), extended by zero
on the \(Q\)-null set, is the derivative when the stated domination
holds. The finite bias counterexample is proved in the boundary item.
\end{proof}
\end{annexitem}

\begin{annexitem}{annex:p15.harmonic-evidence}{thm:p15-harmonic-evidence}{What posterior reciprocal averages estimate}
\begin{proof}
Since \(P=T/I\), integration of \(H=dG/dT\) gives
\(E_PH=I^{-1}\int H\,dT=1/I\). Squaring gives the stated second
moment, possibly infinite. Choose a finite version outside a \(P\)-null
set. Independent posterior draws therefore produce integrable i.i.d.
weights. The strong law gives \(A_M\to1/I>0\), so eventually the
mean is positive and its inverse tends to \(I\).
Linearity makes \(A_M\) unbiased for \(1/I\); it does not make its
nonlinear inverse unbiased.

The relation \(\Pi\ll T\) holds exactly when \(L>0\) \(\Pi\)-almost
everywhere. Indeed \(T\) vanishes on \(\{L=0\}\); on the positive set,
\(T(A)=0\) implies \(\Pi(A\cap\{L\geq1/j\})=0\) for every positive
integer \(j\), hence \(\Pi(A\cap\{L>0\})=0\). Under this relation,
\(d\Pi/dT=1/L\), and
\(E_P(1/L)^2=I^{-1}\int L^{-1}d\Pi\).
If instead \(s=\Pi(L>0)<1\), posterior draws still lie in \(\{L>0\}\)
almost surely, but
\[
 E_P(1/L)=\frac1I\int_{\{L>0\}}d\Pi=s/I.
\]
Here \(s>0\) because \(I>0\); the strong law gives the stated wrong
reciprocal limit \(I/s\). With densities \(g,h\), domination requires
\(g=0\) almost everywhere on \(h=0\), and \(H=g/h\) on \(h>0\).
Substitution gives \(\int H^2dT=\int g^2/h\), proving the explicit
variance requirement.
\end{proof}
\end{annexitem}

\begin{annexitem}{annex:p15.evidence-mc-clt}{thm:p15-evidence-mc-clt}{The i.i.d. central limit theorem and denominator limits}
\begin{proof}
For an ordinary mean, apply \cref{thm:p15-clt-iid} to the integrable
i.i.d. weights with finite variance, including its zero-variance case.
For self-normalized importance put \(U_i=W_i(L_i-I)\). Its mean is
\(I-I=0\), and the assumed second moment is \(v_{\rm SN}\).
Where \(D_M=M^{-1}\sum_iW_i>0\),
\[
 \sqrt M(\widetilde I_M-I)=\frac{M^{-1/2}\sum_iU_i}{D_M}.
\]
The numerator has the asserted normal limit and \(D_M\to1\) almost
surely. The random-denominator result \cref{thm:p15-clt-perturbation}
gives the limit without assuming independence between numerator and
denominator. For harmonic means apply the same CLT to
\(H_i-1/I\), and use the exact identity
\[
 \sqrt M(A_M^{-1}-I)=-\frac{I}{A_M}\sqrt M(A_M-1/I).
\]
Here \(I/A_M\to I^2\) in probability, so the same perturbation theorem
gives variance \(I^4v_H\). Positive limiting denominators imply that
zero-denominator events have probability tending to zero; a fixed default
on those events does not change the weak limits.
\end{proof}
\end{annexitem}

\begin{annexitem}{annex:p15.mc-evidence-boundaries}{thm:p15-mc-evidence-boundaries}{Support, reciprocal bias, infinite variance and dependence}
\begin{proof}
For the first example, direct likelihood draws are Bernoulli with mean
\(\varepsilon\) and variance \(\varepsilon(1-\varepsilon)\).
Dividing the variance of their average by \(I^2\) gives
\((1-\varepsilon)/(M\varepsilon)\). Independence gives probability
\((1-\varepsilon)^M\) of all zero draws. The posterior is the first
point mass, where \(L=1\), so its harmonic estimate equals one.
The target measure \(T\) has mass \(\varepsilon\) on that point;
relative to \(Q\), its derivative is constant \(\varepsilon\), giving
ordinary importance with zero variance. However \(\Pi\not\ll Q\);
the attempted prior weight on the sampled point is \(\varepsilon\),
and self-normalization divides \(\varepsilon\) by itself.

In the second example \(I=(1+3)/2=2\). Proposal weights are \(2/3,2\),
and weighted likelihoods \(2/3,6\). Thus their expectation is 2 and
their second moment is
\((3/4)(4/9)+(1/4)36=28/3\), giving variance \(16/3\).
For one draw self-normalization returns \(L\), whose proposal mean is
\((3/4)1+(1/4)3=3/2\). The posterior is \((1/4,3/4)\);
its one-draw harmonic estimate also returns \(L\), now with expectation
\(1/4+9/4=5/2\). Its reciprocal has expectation
\(1/4+(3/4)(1/3)=1/2\).
To realize these likelihood values with normalized data laws, take the
two data models uniform on \([0,1]\) and \([0,1/3]\), respectively,
and observe an interior point common to both supports.

For the third example the normal density ratio is
\(W(\theta)=\tfrac12 e^{3\theta^2/2}\).
Both densities have full support and \(E_QW=1\), but \(W^2q\)
is a positive constant times \(e^{\theta^2}\), whose integral diverges.
A likelihood identically one can be the probability mass of a one-point
observation space, so this is a proper evidence problem.

For the fourth example, \(\int L^{-1}d\Pi\) is a positive constant times
\[
 \int_{\mathbb R}
 \exp\!\left\{\left(\frac1{2s^2}-\frac1{2\tau^2}\right)\theta^2
                       -\frac{x}{s^2}\theta+\frac{x^2}{2s^2}\right\}d\theta.
\]
If \(\tau^2>s^2\), the positive quadratic term makes it divergent.
At equality, a nonzero linear term diverges along one tail, while a zero
linear term leaves a positive constant over the line. For \(\tau^2<s^2\)
square completion gives a finite Gaussian integral. Evidence itself is
the positive finite normal convolution, so the reciprocal first moment
is finite in every case; only its second moment fails at the stated
threshold.

Finally, the identity transition kernel keeps its initialized posterior
state fixed and is stationary. The harmonic average is then \(1/L\) of
that same state at every simulation count, with inverse \(1/4\) or \(3/4\).
The true evidence is \((1/4+3/4)/2=1/2\). This proves failure of
convergence despite correct posterior marginals and identifies the missing
averaging assumption, without asserting a general MCMC theorem.
\end{proof}
\end{annexitem}

\begin{annexitem}{annex:p15-w02b-p15-056-infobound}{thm:p15-w02b-p15-056-information-replacement-boundary}{Logistic data with observed and expected information}
\begin{proof}
\emph{Logistic facts.} Here \(\log f(u)=-u-2\log(1+e^{-u})\), so
\((\log f)'(u)=-\tanh(u/2)\). Since
\(\operatorname{sech}^2(u/2)=4e^{-u}/(1+e^{-u})^2\),
\[
 (\log f)''(u)=-\tfrac12\operatorname{sech}^2(u/2)=-2f(u)<0 .
\] Also
\(f(u)=\tfrac14\operatorname{sech}^2(u/2)\) is even and decreasing in
\(|u|\), with \(f(0)=1/4\) and \(f(u)\to0\) as \(|u|\to\infty\). Every
derivative of \(\tanh\) is a polynomial in \(\tanh\), so every derivative
of \(\log f\) is bounded on \(\mathbb R\). For a location family the
information per observation at \(\theta\) is
\(E[-(\log f)''(X-\theta)]=2\int f^2\). The substitution
\(t=1/(1+e^{-u})\) has \(dt=f(u)\,du\) and \(f(u)=t(1-t)\), so
\(\int f^2\,du=\int_0^1t(1-t)\,dt=1/6\). The expected information per
observation is therefore \(1/3\) at every \(\theta\).

\emph{First item.} For \(m\) observations at each of \(\pm a\),
\(h_n(\theta)=\tfrac12\{\log f(a-\theta)+\log f(-a-\theta)\}\). Each
summand has second derivative \(-2f<0\), so \(h_n\) is strictly concave.
Also \(h_n'(\theta)=\tfrac12\{\tanh((a-\theta)/2)+\tanh((-a-\theta)/2)\}\)
vanishes at \(\theta=0\) because \(\tanh\) is odd. Hence
\(\widehat\theta_n=0\) is the unique maximizer, and
\(J_n=-h_n''(0)=f(a)+f(-a)=2f(a)\).

\emph{The Laplace hypotheses for \(0\leq a\leq\bar a\).} Take
\(\rho=1\), \(\Theta_n=\mathbb R\), and amplitude the standard normal
density \(\varphi\). The gradient vanishes at zero, and
\(J_n=2f(a)\in[2f(\bar a),1/2]\). The third and fourth derivatives of
\(h_n\) are averages of derivatives of \(\log f\), hence bounded. The
derivatives \(\varphi'\) and \(\varphi''\) are bounded, and
\(\varphi(0)=(2\pi)^{-1/2}\). For \(|v|\leq1\),
\(-h_n''(v)=f(a-v)+f(a+v)\geq2f(\bar a+1)\), because
\(|a\pm v|\leq\bar a+1\). Taylor's theorem with \(h_n'(0)=0\) gives
\(h_n(v)-h_n(0)\leq-c_0v^2\) with \(c_0=f(\bar a+1)\). For
\(|\theta|\geq1\), concavity of \(t\mapsto h_n(t\theta/|\theta|)\)
gives
\(h_n(\theta)-h_n(0)\leq|\theta|\{h_n(\theta/|\theta|)-h_n(0)\}\leq-c_0|\theta|\).
Hence
\[
 \int_{|\theta|\geq1}e^{n(h_n(\theta)-h_n(0))}\varphi(\theta)\,d\theta
 \leq\frac2{\sqrt{2\pi}}\int_1^\infty e^{-nc_0t}\,dt
 \leq\frac{2}{\sqrt{2\pi}\,c_0}e^{-c_0n}.
\]
All the conditions of \cref{def:p15-laplace-family} hold with constants
depending only on \(\bar a\), so \cref{thm:p15-laplace-mode-variants}
gives \(I_n/\widehat I(F_n)=1+O(n^{-1})\). In the approximation only the
determinant changes when \(F_n\) is replaced by \(\widetilde F_n\), so
\(\widehat I(\widetilde F_n)/\widehat I(F_n)=(J_n/\widetilde J_n)^{1/2}\).

\emph{Second item.} The equation \(f(a)=1/6+n^{-1/2}/2\) has a solution
in \([0,\infty)\) exactly when \(1/6+n^{-1/2}/2\leq1/4\), that is when
\(n\geq36\), and the solution is unique by monotonicity. Since
\(f(a_n)>1/6=f(a^*)\) with \(a^*=\log(2+\sqrt3)\), it satisfies
\(a_n<a^*\). Here \(e^{-a^*}=2-\sqrt3\) is the root below one of
\(t^2-4t+1=0\), which is the equation \(t/(1+t)^2=1/6\). Take
\(\bar a=a^*\). Then \(J_n=1/3+n^{-1/2}\), and for \(n\geq38\),
\(\delta_n=n^{-1/2}\leq38^{-1/2}<1/6=c/2\). Thus
\[
 \log\frac{\widehat I(\widetilde F_n)}{I_n}
 =\tfrac12\log(3J_n)-\log\frac{I_n}{\widehat I(F_n)}
 =\tfrac12\log(1+3n^{-1/2})+O(n^{-1}).
\]
Write \(x=n^{-1/2}\). Then
\((1+3x)^{1/2}-1=3x/\{(1+3x)^{1/2}+1\}\), so
\(\sqrt n\{(1+3x)^{1/2}-1\}\to3/2\). Multiplying by \(1+O(n^{-1})\)
changes the scaled quantity by \(O(n^{-1/2})\), which proves the stated
limit.

\emph{Third item.} With \(t=e^{-a}\), \(f(a)=1/24\) is the equation
\(t^2-22t+1=0\). Its root in \((0,1)\) is \(t=11-\sqrt{120}\), so
\(a=-\log(11-\sqrt{120})=\log(11+\sqrt{120})\), because the product of
the two roots is one. Then \(J_n=1/12\) and
\(\|\widetilde J_n-J_n\|=1/3-1/12=1/4\). Any eigenvalue bound \(c\)
satisfies \(c\leq1/12\), so \(c/2\leq1/24<1/4\). With \(\bar a=a\), the
Laplace hypotheses hold as shown above, and
\(\widehat I(\widetilde F_n)/I_n=(1/4)^{1/2}(1+O(n^{-1}))\to1/2\).
The logistic density is positive everywhere, so these observations are
possible under every \(\theta\).
\end{proof}
\end{annexitem}

\begin{annexitem}{annex:p15-w02b-p15-056-newton}{thm:p15-w02b-p15-056-one-step-newton}{Expansion identity and error of the one-step formula}
\begin{proof}
\emph{First item.} Put \(u=\bar\theta-\widehat\theta=(F+G)^{-1}\lambda'\),
writing \(\lambda'=\lambda'(\widehat\theta)\). Then
\(\ell^{(2)}(\bar\theta)=\ell(\widehat\theta)-\tfrac12u^TFu\) and
\(\lambda^{(1)}(\bar\theta)=\lambda(\widehat\theta)+\lambda'^Tu\).
Because \((F+G)^{-1}\) is symmetric,
\[
 2\lambda'^Tu-u^TFu
 =\lambda'^T(F+G)^{-1}\{2-F(F+G)^{-1}\}\lambda' .
\]
Adding \(2\ell(\widehat\theta)+2\lambda(\widehat\theta)+d\log(2\pi)-\log|F+G|\)
gives \(2\ell(\widehat\theta)+E\), as claimed. Formula (4) is
\((2\pi)^{d/2}|\widetilde\Sigma|^{1/2}\) times the likelihood and prior
at the posterior mode. Twice its logarithm, with the three stated
replacements, is the right side.

\emph{Second item.} Here \(\lambda'=D\pi/\pi\) at the mode. The amplitude
bounds give \(\|\lambda'\|\leq A/a_-\). The matrix \(G=W^{-1}\) is
positive definite, so \(F+G\succeq F\succeq ncI\). Hence
\(\|(F+G)^{-1}\|\leq1/(nc)\). Since \(\|F\|\leq nC\),
\(\|2I-F(F+G)^{-1}\|\leq2+C/c\). The quadratic term of \(E\) is
therefore at most \((A/a_-)^2(2+C/c)/(nc)\) in absolute value. Next,
\(\log|F+G|-\log|F|=\log\det(I+M)\) with
\(M=F^{-1/2}GF^{-1/2}\) symmetric, positive semidefinite and of norm at
most \(g/(nc)\). Its eigenvalues \(\mu_j\) lie in \([0,g/(nc)]\), and
\(0\leq\log(1+\mu_j)\leq\mu_j\), so this difference lies in
\([0,dg/(nc)]\). With
\(2\log\widehat I_{\rm MLE}=2\ell(\widehat\theta)+2\lambda(\widehat\theta)+d\log(2\pi)-\log|F|\),
the two bounds give
\(2\ell(\widehat\theta)+E-2\log\widehat I_{\rm MLE}=O(n^{-1})\).
\Cref{thm:p15-laplace-mode-variants} gives
\(2\log I_n-2\log\widehat I_{\rm MLE}=O(n^{-1})\). Adding the two
proves the claim for one model. Subtracting the claims for the two
models, \(2\ell_1(\widehat\theta_1)-2\ell_0(\widehat\theta_0)=\Lambda\),
gives the factor statement.

\emph{Third item.} Here \(\ell_n(\theta)=\log\phi_{\theta,s^2}(x_n)\)
is quadratic, with \(\widehat\theta=x_n\) and \(F=1/s^2\). The prior is
\(\lambda=\log\phi_{\mu,\tau^2}\), with \(G=1/\tau^2\) and
\(\lambda'(x_n)=-D_n/\tau^2\). Since \(\ell=\ell^{(2)}\) and
\(\lambda=\lambda^{(1)}-\tfrac12G(\theta-\widehat\theta)^2\), the
first item gives
\[
 \ell(\widehat\theta)+\tfrac12E
 =\ell(\bar\theta)+\lambda(\bar\theta)+\tfrac12Gu^2
 +\tfrac12\log(2\pi)-\tfrac12\log(F+G).
\]
The log posterior kernel \(\ell+\lambda\) is quadratic with curvature
\(F+G\), so one Newton step from any point lands at its maximizer, and
\(\bar\theta=\widetilde\theta_n\). The posterior-mode calculation is exact
by \cref{thm:p15-laplace-normal-example}, so the terms other than
\(\tfrac12Gu^2\) sum to \(\log I_n\). Finally
\(u=-(D_n/\tau^2)s^2\tau^2/(s^2+\tau^2)=-D_ns^2/(s^2+\tau^2)\), and
\(\tfrac12Gu^2=D_n^2s^4/\{2\tau^2(s^2+\tau^2)^2\}\). Since
\(s^4=\sigma^4/n^2\), this is \(O(n^{-2})\) for bounded \(D_n\). The
expression for approximation (5) is the logarithm of the ratio in
\cref{thm:p15-laplace-normal-example}. At \(D_n=0\) it equals
\(\tfrac12\log\{1+\sigma^2/(n\tau^2)\}\), and
\(n\) times this tends to \(\sigma^2/(2\tau^2)>0\).
\end{proof}
\end{annexitem}

\begin{annexitem}{annex:p15-w02b-p15-056-unitinfo}{thm:p15-w02b-p15-056-unit-information-schwarz}{The constant-order terms under the unit-information prior}
\begin{proof}
By the first premise,
\(\log\pi_1(\widehat\beta,\widehat\psi)=\log\pi_0(\widehat\beta)
 -\tfrac p2\log(2\pi)-\tfrac12\log\det W
 -\tfrac12(\widehat\psi-\psi_0)^TW^{-1}(\widehat\psi-\psi_0)\).
Substituting this into the definitions of \(c_{1,n}\) (dimension
\(d_0+p\)) and \(c_{0,n}\) (dimension \(d_0\)) gives
\[
\begin{split}
 c_{1,n}-c_{0,n}={}&\{\log\pi_0(\widehat\beta_n)-\log\pi_0(\widehat\beta^*_n)\}
 -\tfrac12(\widehat\psi_n-\psi_0)^TW^{-1}(\widehat\psi_n-\psi_0)\\
 &-\tfrac12\{\log\det J_{1,n}-\log\det J_{0,n}+\log\det W\}.
\end{split}
\]
By the second premise, the first brace is at most \(K^2n^{-1/2}\) in
absolute value, and the quadratic term is at most
\(\tfrac12\|W^{-1}\|K^2n^{-1}\).

For the determinants, \(A_n\) and \(D_n\) are principal blocks of
\(J_{1,n}\). Their Rayleigh quotients are Rayleigh quotients of
\(J_{1,n}\) on coordinate subspaces, so their eigenvalues lie in
\([c,C]\). Write \(\Sigma_n=D_n-B_n^TA_n^{-1}B_n\). The factorization
\[
 J_{1,n}=
 \begin{pmatrix}I&0\\B_n^TA_n^{-1}&I\end{pmatrix}
 \begin{pmatrix}A_n&0\\0&\Sigma_n\end{pmatrix}
 \begin{pmatrix}I&A_n^{-1}B_n\\0&I\end{pmatrix}
\]
can be checked by multiplication. Its outer factors are unit triangular,
so \(\det J_{1,n}=\det A_n\det\Sigma_n\). The proof of
\cref{thm:p15-laplace-curvature-replacement} shows the following. If a
symmetric \(J\) has eigenvalues at least \(c\) and \(\|\widetilde J-J\|\leq\delta<c/2\),
then \(|\log\det\widetilde J-\log\det J|\leq2\dim(J)\delta/c\). Apply it
twice. First take \(J=J_{0,n}\) and \(\widetilde J=A_n\), with
\(\delta=Kn^{-1/2}\). Then take \(J=D_n\) and \(\widetilde J=\Sigma_n\),
with \(\delta=\|B_n^TA_n^{-1}B_n\|\leq K^2/(cn)\). For large \(n\) both
values of \(\delta\) are below \(c/2\), and
\[
\begin{split}
 &\log\det J_{1,n}-\log\det J_{0,n}+\log\det W\\
 &\quad=\{\log\det A_n-\log\det J_{0,n}\}\\
 &\qquad+\{\log\det\Sigma_n-\log\det D_n\}+\{\log\det D_n+\log\det W\}.
\end{split}
\]
The three braces are \(O(n^{-1/2})\), \(O(n^{-1})\) and, by the third
premise, \(O(n^{-1/2})\). Hence \(c_{1,n}-c_{0,n}=O(n^{-1/2})\). Since
\(\mathrm{BIC}_k=-2\ell_k+d_k\log n\) and \(d_1-d_0=p\),
\(-\tfrac12(\mathrm{BIC}_1-\mathrm{BIC}_0)=S_n\).
\Cref{thm:p15-bic-leading-order} gives
\(\log B_{10}=S_n+c_{1,n}-c_{0,n}+O(n^{-1})=S_n+O(n^{-1/2})\).

For the normal instance apply \cref{thm:p15-normal-point-null-factor}
with \(\theta_0=\mu_1=0\), \(s^2=\sigma^2/n\) and \(\tau^2=\sigma^2\),
so that \(r=n\) and \(z=\bar x/s\). This gives
\(\log B_{10}=-\tfrac12\log(1+n)+nz^2/\{2(n+1)\}\). The full-sample
log-likelihood ratio at the estimate \(\widehat\psi=\bar x\) is
\(n\bar x^2/(2\sigma^2)=z^2/2\), so \(S_n=z^2/2-\tfrac12\log n\).
Subtracting gives the displayed difference, which is at most
\((1+K^2)/(2n)\) in absolute value when \(|z|\leq K\). This instance
meets the premises: \(d_0=0\), \(J_{1,n}=D_n=1/\sigma^2=W^{-1}\), and
\(|\widehat\psi_n|=\sigma|z|n^{-1/2}\). The hypotheses of
\cref{thm:p15-bic-leading-order} for bounded \(\bar x\) were verified in
\cref{thm:p15-laplace-normal-example}. If \(\bar x\to\psi_1\neq0\), then
\(z^2/\{2(n+1)\}=n\bar x^2/\{2\sigma^2(n+1)\}\to\psi_1^2/(2\sigma^2)\),
and the logarithmic term tends to zero.
\end{proof}
\end{annexitem}

\begin{annexitem}{annex:p15-w02b-p15-056-mixture}{thm:p15-w02b-p15-056-mixture-estimators}{Monotone estimating equations}
\begin{proof}
\emph{Shape of \(\psi_p\).} For \(p>0\) the denominator of \(\psi_p\)
is positive. Differentiation gives
\[
 \partial_p\psi_p=\frac{L}{\{\delta p+(1-\delta)L\}^2}\geq0,
 \qquad
 \partial_p^2\psi_p=-\frac{2\delta L}{\{\delta p+(1-\delta)L\}^3}.
\]
If \(L=0\), then \(\psi_p=1/\delta\) for all \(p\). If \(L>0\), then
\(\psi_p\to-1/(1-\delta)\) as \(p\downarrow0\) and \(\psi_p\to1/\delta\)
as \(p\to\infty\), strictly increasing. Since \(p/(\delta p+(1-\delta)L)\leq1/\delta\)
and \(L/(\delta p+(1-\delta)L)\leq1/(1-\delta)\),
\(|\psi_p|\leq\max\{1/\delta,1/(1-\delta)\}\). With \(a=\delta p\) and
\(b=(1-\delta)L\), the inequalities \((a+b)^2\geq4ab\) and
\((a+b)^3\geq3a^2b\) give
\(0\leq\partial_p\psi_p\leq1/\{4\delta(1-\delta)p\}\) and
\(|\partial_p^2\psi_p|\leq2/\{3\delta(1-\delta)p^2\}\).

\emph{First item.} The stated limits and monotonicity of \(\Psi_m\)
follow by summation. If \(N_+=0\), then \(\Psi_m\equiv m/\delta>0\) has
no root. If \(N_+\geq1\), \(\Psi_m\) is continuous and strictly
increasing, so the intermediate value theorem gives a unique root exactly
when \(\Psi_m(0+)<0\). For the \(\widehat{\rm pr}_4\) equation, write
\(N(p)\) and \(D(p)\) for its numerator and denominator. Then \(p=N(p)/D(p)\)
if and only if \(pD(p)=N(p)\), that is
\(\sum_iw_p(\theta_i)\{p-L(\theta_i)\}=0\). The imagined terms cancel.
Since \(w_p(p-L)=p\,\psi_p\), this is \(\sum_i\psi_p(\theta_i)=0\). Each
posterior draw lies in \(\{L>0\}\) almost surely, because
\(P(L=0)=I^{-1}\int_{\{L=0\}}L\,d\Pi=0\). So \(N_0=0\), \(N_+=m\), and
the unique root exists.

\emph{Second item.} The measure \(Q_\delta\) has density
\(\delta+(1-\delta)L/I\) relative to \(\Pi\). Let
\(g(p)=E_{Q_\delta}\psi_p\). At \(p=I\),
\[
 g(I)=\int\frac{I-L}{\delta I+(1-\delta)L}\cdot
 \frac{\delta I+(1-\delta)L}{I}\,d\Pi=\frac{I-I}{I}=0 .
\]
Bounded convergence gives \(g'(p)=E_{Q_\delta}\partial_p\psi_p\). This
is positive, because \(Q_\delta(L>0)\geq(1-\delta)P(L>0)=1-\delta\). So
\(g\) is strictly increasing with unique root \(I\). Fix a rational
\(\epsilon\in(0,I)\). The bounded i.i.d. strong law gives
\(\Psi_m(I\pm\epsilon)/m\to g(I\pm\epsilon)\) almost surely, with
\(g(I-\epsilon)<0<g(I+\epsilon)\), and \(N_+/m\to Q_\delta(L>0)>0\).
Almost surely, therefore, eventually \(N_+\geq1\) and
\(\Psi_m(I-\epsilon)<0<\Psi_m(I+\epsilon)\). Then the root exists, is
unique, and lies in \((I-\epsilon,I+\epsilon)\). Intersecting over
rational \(\epsilon\) proves almost-sure convergence.

For the limit law, on the event that the root \(\widehat p\) exists and
lies in \([I/2,2I]\), the mean value theorem gives \(\bar p\) between
\(\widehat p\) and \(I\) with
\[
 \sqrt m(\widehat p-I)=-\frac{m^{-1/2}\Psi_m(I)}{m^{-1}\Psi_m'(\bar p)} .
\]
The numerator is a normalized sum of bounded i.i.d. variables with mean
\(g(I)=0\). By \cref{thm:p15-clt-iid} it converges to
\(\operatorname{normal}(0,E_{Q_\delta}\psi_I^2)\). On \([I/2,2I]\) the
second-derivative bound makes \(m^{-1}\Psi_m'\) Lipschitz with constant
\(8/\{3\delta(1-\delta)I^2\}\), uniformly in the sample. Hence
\(|m^{-1}\Psi_m'(\bar p)-m^{-1}\Psi_m'(I)|\to0\) almost surely. The strong
law gives \(m^{-1}\Psi_m'(I)\to g'(I)>0\). The complementary event has
probability tending to zero. \Cref{thm:p15-clt-perturbation}, with a fixed
default value on that event, gives the normal limit with variance
\(E_{Q_\delta}\psi_I^2/g'(I)^2=v_3\).

\emph{Third item.} Let \(g_4(p)=E_P\psi_p\). Under \(P\), \(L>0\)
almost surely, so \(g_4\) is continuous and strictly increasing, with
limits \(-1/(1-\delta)\) and \(1/\delta\) by bounded convergence. It
therefore has a unique root \(p_\delta\). The argument of the second
item, with \(P\) in place of \(Q_\delta\), proves that the roots
converge almost surely to \(p_\delta\). Since \(g_4\) is strictly
increasing, \(p_\delta=I\) exactly when \(g_4(I)=0\). Because
\(dP/d\Pi=L/I\), that is the displayed integral condition. In the
example, \(P=(1/4,3/4)\) and, with \(\delta=1/2\),
\[
 g_4(p)=\frac12\cdot\frac{p-1}{p+1}+\frac32\cdot\frac{p-3}{p+3}.
\]
Clearing denominators gives \(2p^2-2p-6=0\), whose positive root is
\((1+\sqrt{13})/2\neq2\). For \(\widehat{\rm pr}_3\),
\(Q_{1/2}=(3/8,5/8)\), and
\(g(2)=\tfrac38\cdot\tfrac{1}{3/2}-\tfrac58\cdot\tfrac{1}{5/2}=0\),
in agreement with the second item.
\end{proof}
\end{annexitem}

\begin{annexitem}{annex:p15-w02b-p15-056-bridge}{thm:p15-w02b-p15-056-bridge-identity}{Two expectations}
\begin{proof}
The function \(\alpha q_1\) is \(p_2\)-integrable because
\[
 \int|\alpha q_1|\,p_2\,d\mu=c_2^{-1}\int|\alpha|q_1q_2\,d\mu<\infty .
\]
Its integral is \(\kappa/c_2\). In the same way \(E_{p_1}[\alpha q_2]=\kappa/c_1\).
Since \(\kappa\neq0\), dividing gives \(c_1/c_2\). The two sample means
are averages of integrable i.i.d. variables. By
\cref{thm:p15-etemadi-strong-law} they converge almost surely to
\(\kappa/c_2\) and \(\kappa/c_1\). The second limit is nonzero, so the
denominator is eventually nonzero and the ratio converges.

For the first special case,
\(E_f[\alpha q_1]=\int_{\{f>0\}}L\pi\,d\mu=I\,P(f>0)\). For the second,
\(E_f[\alpha q_1]=\int_{\{L\pi>0\}}f\,d\mu\) and
\(E_P[\alpha f]=E_P[f/(L\pi)]\). In the first case
\(\int|\alpha|q_1q_2\,d\mu=I\,P(f>0)\leq I\). In the second it is
\(\int_{\{L\pi>0\}}f\,d\mu\leq1\). The first display is the direct
computation. The identity applies to the second whenever the integral is
positive, and rearrangement gives the displayed form. For two models on
a common space, \(c_1=I_1\) and \(c_2=I_0\).
\end{proof}
\end{annexitem}

\begin{annexitem}{annex:p15-w02b-p15-056-lapmet}{thm:p15-w02b-p15-056-laplace-metropolis}{Mode and covariance errors in the Laplace--Metropolis value}
\begin{proof}
Since \(|\Sigma^*_n|=|H_n|^{-1}\det(\mathrm{I}+E_n)\), the logarithms of
\(\widehat I_{\rm LM}\) and of
\(\widehat I_{\rm post}=(2\pi)^{d/2}|H_n|^{-1/2}e^{q_n(\widetilde\theta_n)}\)
differ by the displayed amount. Put
\(v=\theta^*_n-\widetilde\theta_n\). The third premise of
\cref{def:p15-laplace-family} gives
\(q_n(\widetilde\theta_n)-q_n(\theta^*_n)\geq nc_0\|v\|^2\). For the upper
bound, the gradient at the mode is zero, and Taylor's theorem with
integral remainder gives
\[
 q_n(\widetilde\theta_n+v)-q_n(\widetilde\theta_n)
 =n\int_0^1(1-t)\,D^2h_n(\widetilde\theta_n+tv)[v,v]\,dt .
\]
On the ball,
\(\|D^2h_n(\xi)\|\leq\|J_n\|+\sup\|D^3h_n\|\,\|\xi-\widetilde\theta_n\|\leq C+A\rho\).
Hence the absolute difference is at most \(\tfrac12n(C+A\rho)\|v\|^2\).
The eigenvalues of \(E_n\) lie in \([-e_n,e_n]\subset(-1/2,1/2)\), so
\(\mathrm{I}+E_n\) is positive definite. Since \(|\log(1+t)|\leq2|t|\)
there, \(|\tfrac12\log\det(\mathrm{I}+E_n)|\leq de_n\).

\Cref{thm:p15-laplace-mode-variants} gives
\(\log I_n-\log\widehat I_{\rm post}=O(n^{-1})\). Adding the bounds
proves the stated orders and the inheritance condition. If \(e_n\to0\)
and \(\liminf n\varepsilon_n^2>0\), then
\(\log\widehat I_{\rm post}-\log\widehat I_{\rm LM}\geq nc_0\varepsilon_n^2-de_n\).
This has positive \(\liminf\), and so does \(\log I_n-\log\widehat I_{\rm LM}\).

In the normal instance, the log posterior kernel is
\(q(\theta)=q(\widetilde\theta)-\tfrac12H(\theta-\widetilde\theta)^2\)
with \(H=1/s^2+1/\tau^2\), and the posterior is
\(\operatorname{normal}(\widetilde\theta,1/H)\). Exact draws are
\(\theta_i=\widetilde\theta+Z_iH^{-1/2}\). The maximizer of \(q\) among
them minimizes \(|Z_i|\); ties have probability zero. Hence
\(q(\widetilde\theta)-q(\theta^*)=\tfrac12\min_iZ_i^2\), and with
\(\Sigma^*=1/H\), \(E_n=0\). The posterior-mode calculation is exact by
\cref{thm:p15-laplace-normal-example}, so
\(\log I_n-\log\widehat I_{\rm LM}=\tfrac12\min_iZ_i^2\). Independence
gives \(\mathbb P(\tfrac12\min_iZ_i^2>t)=\mathbb P(|Z_1|>\sqrt{2t})^M
=\{2-2\Phi(\sqrt{2t})\}^M\). This law does not involve \(n\), and the
error is positive almost surely because \(\mathbb P(Z_i=0)=0\).
\end{proof}
\end{annexitem}

\begin{annexitem}{annex:p15-w02b-p15-056-sdratio}{thm:p15-w02b-p15-056-savage-dickey}{Evaluating the posterior marginal at the null}
\begin{proof}
By definition of the conditional version,
\[
 p_1(\psi_0\mid D)=\frac{\pi_1(\psi_0)}{I_1}
 \int L(\beta,\psi_0)\pi_1(\beta\mid\psi_0)\,d\beta .
\]
\emph{First item.} Under the condition, the integral is \(I_0\). Thus
\(p_1(\psi_0\mid D)=\pi_1(\psi_0)I_0/I_1>0\), and
\(\pi_1(\psi_0)/p_1(\psi_0\mid D)=I_1/I_0=B_{10}\).

\emph{Second item.} The support hypothesis lets the integral be taken
over \(\{\pi_0>0\}\):
\[
\begin{split}
 \int L(\beta,\psi_0)\pi_1(\beta\mid\psi_0)\,d\beta
 &=\int_{\{\pi_0>0\}}\frac{\pi_1(\beta\mid\psi_0)}{\pi_0(\beta)}
 L(\beta,\psi_0)\pi_0(\beta)\,d\beta\\
 &=I_0E_{P_0}\!\left[\frac{\pi_1(\beta\mid\psi_0)}{\pi_0(\beta)}\right],
\end{split}
\]
because \(P_0(d\beta)=L(\beta,\psi_0)\pi_0(\beta)\,d\beta/I_0\) gives
\(\{\pi_0=0\}\) mass zero. The equality holds in \([0,\infty]\).
Dividing by \(I_1\) gives the identity.

\emph{Third item.} Since \(p\geq1\), the set \(\{\psi=\psi_0\}\) is
Lebesgue-null in \(\mathbb R^{d_0+p}\). Changing \(\pi_1\) there changes
no integral over that space, so \(I_1\) is unchanged. \(I_0\) does not
involve \(\pi_1\). The new marginal value is \(\int g\), and the new
posterior marginal value is \(I_1^{-1}\int L(\beta,\psi_0)g(\beta)\,d\beta\).

\emph{Fourth item.} With \(d_0=0\) the conditional condition is vacuous.
Here \(I_0=L(\psi_0)\) and \(p_1(\psi_0\mid x)=L(\psi_0)\pi_1(\psi_0)/I_1\),
so the ratio is \(I_1/L(\psi_0)=m_1(x)/m_0(x)\), the reciprocal of the
factor in \cref{thm:p15-normal-point-null-factor}. In the two-coordinate
example, \(\pi_1(\beta\mid0)=\phi_{0,1}(\beta)\neq\phi_{0,4}(\beta)=\pi_0(\beta)\),
and both are positive. The posterior under \(H_0\) at \(x_1=0\) is
\(\operatorname{normal}(0,4/5)\), since precisions add:
\(1/4+1=5/4\). Here \(\phi_{0,1}/\phi_{0,4}=2e^{-3\beta^2/8}\). For
\(\beta\sim\operatorname{normal}(0,v)\) and \(a>0\), completing the
Gaussian integral gives \(Ee^{-a\beta^2}=(1+2av)^{-1/2}\). With
\(a=3/8\) and \(v=4/5\) this equals \((5/8)^{1/2}\), so the correction
factor is \(2(5/8)^{1/2}=\sqrt{5/2}\).
\end{proof}
\end{annexitem}

\begin{annexitem}{annex:p15-w02b-p15-056-latent}{thm:p15-w02b-p15-056-latent-ratio}{Latent-variable ratios and their support condition}
\begin{proof}
\emph{First item.}
\[
 E_{P^z_0}B_{10}(z)=\int_{\{p_0>0\}}\frac{p_1}{p_0}\cdot\frac{p_0}{m_0}\,d\nu
 =\frac1{m_0}\int_{\{p_0>0\}}p_1\,d\nu=\frac{m_1}{m_0}P^z_1(p_0>0).
\]
Because \(p_1\geq0\), \(P^z_1(p_0>0)=1\) exactly when
\(\int_{\{p_0=0\}}p_1\,d\nu=0\).

\emph{Second item.} The ratios are nonnegative i.i.d. variables with the
finite mean just computed, so \cref{thm:p15-etemadi-strong-law} applies.
Their second moment is
\[
 m_0^{-1}\int_{\{p_0>0\}}p_1^2/p_0\,d\nu,
\]
which gives the variance criterion.

\emph{Third item.} Apply the first two items with the models
exchanged. Then \(p_0/p_1=L_0/L_1\) on \(\{p_1>0\}\), the draws come from
\(P^z_1\), the posterior of \(\theta_1\), and the limit is
\(B_{01}P^z_0(p_1>0)\). Here \(P^z_0\) has density \(p_0/m_0=\pi_1\),
since \(m_0=L_0\). Also \(\{p_1>0\}=\{L_1>0,\pi_1>0\}\), so
\(P^z_0(p_1>0)=\pi_1(L_1>0)\). The average of \(1/L_1(z_i)\) is the
statistic \(A_M\) of \cref{thm:p15-harmonic-evidence} with \(G=\Pi\).
\end{proof}
\end{annexitem}
